\chapter{Conclusions and Future Work}
\label{chap:conclusions}

\section{Conclusions}

This dissertation investigated whether a reproducible \gls{xai} framework could be applied across heterogeneous replications of medical-prediction studies and what its outputs could validly reveal about model behaviour. Five published workflows concerning heart disease, stroke, diabetes, \gls{ckd}, and \gls{hbaone}-associated hospital readmission were reconstructed. One final analytical configuration from each replication was then examined using native importance where available, permutation importance, \gls{shap}, representative-case \gls{lime}, rank agreement, and bootstrap-refit stability where retraining was possible.

The work produced a documented replication-and-correction process and a reusable framework for generating and comparing explanation artefacts across different models and feature representations. The results demonstrate that explanation quality depends partly on the validity and reproducibility of the underlying predictive workflow. \gls{xai} cannot correct information leakage, ambiguous preprocessing, test-set-driven model selection, or an evaluation design that does not support the intended claim \citep{collins_tripodai_2024, moons_probastai_2025}.

\subsection{Answer to RQ1}

All five source workflows could be reconstructed at the level of dataset, principal preprocessing stages, model families, and analytical objective. Exact numerical reproduction was not always possible because some procedures depended on incompletely reported random choices, implementation-specific neural-network details, ambiguous preprocessing order, or differences between predictive and inferential designs.

Separating paper-style replication from corrected analysis preserved the original procedures while preventing known methodological limitations from entering the final \gls{xai} configurations. In \gls{ckd}, the two branches exposed the effect of fitting preprocessing before the data partition. In Stroke and Diabetes, learned preprocessing and feature selection were restricted to the training data. In Heart, logistic regression was selected from the replicated analytical design before the held-out ranking was examined. This prevented the higher observed \gls{rocauc} of the backpropagation model from being used retrospectively to change the final \gls{xai} configuration. \gls{hbaone} retained its full-cohort inferential design and was not presented as an out-of-sample classifier.

RQ1 is answered with a qualified conclusion. The workflows were procedurally reproducible, while exact numerical equivalence was constrained by reporting detail and implementation choices. Explicit assumptions and corrected analytical branches produced defensible final configurations without obscuring how they differed from the source workflows.

\subsection{Answer to RQ2}

Agreement among global explanation methods depended on the dataset and model representation. Heart, Stroke, and \gls{ckd} showed strong agreement between native importance and \gls{shap} but lower agreement involving permutation importance. Diabetes produced the same ordering under all three methods, although its five-feature design made complete top-five membership inevitable. In \gls{hbaone}, the absolute-coefficient ranking differed markedly from the \gls{shap} ranking. Permutation importance and \gls{shap} were instead closely aligned.

Native importance, permutation importance, and \gls{shap} are therefore complementary rather than interchangeable. They quantify different aspects of a fitted model and respond differently to encoding, sparsity, interactions, predictor dependence, and estimator structure \citep{fisher_models_2019, NIPS2017_8a20a862, vilone_classification_2021}. Agreement does not prove clinical correctness, and disagreement does not by itself indicate that a method has failed. RQ2 is answered by showing that consistency must be evaluated empirically for each fitted model.

\subsection{Answer to RQ3}

Repeated \gls{lime} explanations were generated for up to ten observations per available correct, incorrect, and borderline category under 20 perturbation seeds. Mean seed-to-seed rank correlations ranged from 0.898 to 0.969, showing that the leading local rules were generally persistent. The same leading rule nevertheless appeared in both correct and incorrect illustrative cases. High age supported a Stroke true positive and a high-confidence false positive, while low age supported a true negative and a false negative. Similar patterns occurred for chest-pain type in Heart, glucose in Diabetes, and serum creatinine or haemoglobin in \gls{ckd}.

Global-local feature presence was high in several datasets, but this measure recorded only whether the same feature names appeared in the global \gls{shap} ranking and local \gls{lime} rules. It did not compare attribution direction, magnitude, \gls{lime} intervals, or prediction correctness. Complete overlap in Diabetes was also partly structural because the model contained only five predictors. \gls{hbaone} showed low global-local feature presence and produced local rules based on encoded interactions, which were less directly interpretable than raw clinical measurements.

RQ3 is answered descriptively with reduced dependence on a single case or random perturbation sequence. \gls{lime} clarified how local surrogates represented selected predictions, including errors, but did not independently determine whether the model reasoning or prediction was correct \citep{ribeiro_why_2016, vilone_classification_2021}. Representative explanations should therefore be presented with the observed outcome, predicted class, prediction probability, case-selection rationale, and sensitivity summary.

\subsection{Answer to RQ4}

Bootstrap-refit stability was evaluated using 100 stratified resamples of each final training matrix, a fixed held-out test set, and negative-log-loss permutation importance. Mean complete-ranking correlations were 0.700 for Heart, 0.812 for Stroke, 0.849 for Diabetes, and 0.949 for \gls{ckd}. Diabetes had the highest top-three overlap, while \gls{ckd} had the highest complete-ranking correlation. Heart showed the lowest persistence under changes in training-sample composition, although \emph{ca} and \emph{thal} remained its most frequent leading features.

The revised \gls{ckd} result demonstrates the importance of the scoring rule. Replacing accuracy with negative log loss removed the coarse score increments and numerous ties that had made the earlier ranking unstable. The final results do not support a simple relationship between predictive performance, within-fit method agreement, and stability after refitting.

No \gls{hbaone} estimate was included in the principal stability comparison because its saved analytical bundle did not contain a validated retraining callback capable of generating equivalent bootstrap refits. A substitute model would have evaluated a different analytical configuration.

An exploratory grouped-permutation bootstrap run showed feasible computation but only 58 of 100 refits converged. Because excluding 42 failed refits could introduce selection bias, its correlation estimate was not adopted as the official \gls{hbaone} stability result. Stability for this configuration therefore remains unresolved.

RQ4 is answered within the implemented bootstrap-refit design. Global rankings showed dataset-dependent persistence under changes in the composition of the final training matrix, and complete-ranking correlation did not always follow the same ordering as top-three overlap.

\subsection{Answer to RQ5}

The framework operated successfully across the five tested configurations. It supported the represented linear classifiers, tree ensembles, Statsmodels logistic-regression wrapper, DataFrame inputs, and encoded design matrix. All configurations produced native importance, permutation importance, \gls{shap}, representative cases, and \gls{lime} explanations. Bootstrap-refit stability was produced for every configuration that supplied the required retraining procedure.

The scientific scope was narrower than the technical coverage. The represented model-family groups were small, and differences in model family coincided with differences in dataset, feature count, sample size, class distribution, and analytical purpose. Measured bundle reruns further showed that computational cost was configuration-dependent: repeated-seed \gls{lime} dominated the four held-out predictive configurations, while model-agnostic \gls{shap} over the large \gls{hbaone} design matrix raised its total execution time to 36.9 minutes. \gls{hbaone} therefore demonstrated compatibility with a substantially different representation while also exposing limitations in local interpretability, computation, held-out evaluation, and retraining.

RQ5 is therefore answered within the tested scope. The framework architecture transferred across the evaluated binary-classification interfaces, while interpretation remained dependent on study design, feature representation, and retraining availability. This distinction prevents demonstrated software compatibility from being overstated as universal scientific generalisability.


\section{Principal Contributions}

The dissertation makes five principal contributions.

\begin{enumerate}
    \item It provides a transparent reconstruction of five heterogeneous medical-data studies and documents ambiguities, leakage risks, class imbalance, and differences between predictive and inferential designs.

    \item It separates source-study reproduction from corrected analytical configurations, preventing information leakage and test-set-driven model selection from being incorporated silently into the \gls{xai} extension.

    \item It provides a reusable binary-classification framework integrating native importance, permutation importance, \gls{shap}, \gls{lime}, representative-case selection, rank comparison, and optional bootstrap-refit stability.

    \item It demonstrates empirically that agreement among explanation methods, local rule behaviour, and ranking stability can differ substantially across datasets and model representations. Strong predictive performance and within-fit agreement did not guarantee stability after retraining.

    \item It provides a reproducibility workflow comprising dependency locks, configuration files, saved analytical bundles, runtime provenance, automated tests, command-line execution, timestamped reruns, and tolerance-aware comparison of regenerated results.
\end{enumerate}

The primary contribution is methodological and computational. The dissertation does not introduce a new attribution algorithm or claim to validate a clinical decision-support system.

\section{Limitations}

The conclusions are bounded by the limitations discussed in Chapter~\ref{chap:discussion}. The datasets were retrospective, heterogeneous, and not externally validated. One final model was explained per dataset, and alternative defensible models could produce different results. Four tasks used held-out prediction, while \gls{hbaone} used full-cohort inference and cannot be treated as directly equivalent.

The framework was evaluated only for binary classification and the model interfaces represented in the experiments. Explanation methods retained their known dependence on data representation, reference distributions, feature dependence, perturbation procedures, scoring metrics, and random variation
\citep{NIPS2017_8a20a862, ribeiro_why_2016,
vilone_classification_2021, lofstrom_meta_2022}. Rank agreement and top-feature overlap assessed similarity between method outputs, not clinical truth. Bootstrap-refit stability used 100 resamples of each final aligned training matrix, a fixed held-out test set, and one global permutation-ranking definition based on negative log loss. It did not characterise every source of explanation uncertainty. Representative local cases were deliberately selected and were not a population sample.

No prospective data, external hospital cohort, temporal validation, fairness audit, causal analysis, clinician study, usability evaluation, calibration intervention, clinical-utility assessment, or regulatory conformity assessment was conducted. The models and explanations must therefore not be interpreted as ready for patient-care use.

These limitations are consistent with contemporary guidance requiring prediction models to be assessed for transparent reporting, risk of bias, applicability, and clinical evaluation before deployment \citep{moons_probastai_2025, collins_tripodai_2024, vasey_decideai_2022}.

\section{Future Work}

\subsection{External and Temporal Validation}

The most important next step is to evaluate the final analytical configurations on independent cohorts. External evaluation would test whether predictive performance and explanation rankings persist under differences in population, clinical practice, coding, and measurement.
Where longitudinal data are available, temporal evaluation would provide evidence about performance on later observations. These extensions should follow established guidance for transparent reporting, risk-of-bias assessment, and applicability of clinical prediction models \citep{collins_tripodai_2024, moons_probastai_2025}.

Model calibration should also be examined explicitly. This is particularly important for imbalanced outcomes such as Stroke and \gls{hbaone}, where discrimination, threshold behaviour, and predicted-probability reliability address different questions. Calibration plots, calibration intercepts, and calibration slopes
should accompany discrimination measures where appropriate \citep{van_calster_calibration_2019}. Thresholds should be selected according to a clearly defined use case and associated error costs rather than adopted automatically at 0.5.

\subsection{Broader Stability Analysis}

Future experiments should separate multiple sources of explanation variability. These may include resampling of the training cohort, alternative train--test partitions, model hyperparameter uncertainty, \gls{shap} background selection, repeated \gls{lime} perturbation, correlated-feature grouping, and controlled perturbations of individual observations. Increasing the number of repeated fits would allow uncertainty intervals for rank and top-feature frequency to be estimated more reliably.

For \gls{hbaone}, future work should use a validated retraining callback with a fitting strategy that converges reliably across resamples and should prespecify how failed fits will be handled. The exploratory grouped-permutation run demonstrated feasible runtime but insufficient convergence (58 of 100 refits), so the successful subset should not be treated as an unbiased stability sample. A simplified substitute model should not be used unless it is introduced as a separate analytical experiment.

\subsection{Model and Task Coverage}

The framework could be extended to multiclass classification, regression, survival analysis, and calibrated risk prediction. Each extension would require task-specific metrics, representative-case definitions, output handling, and tests rather than a superficial relaxation of the current binary checks.

Additional estimators could include calibrated boosting models, interpretable generalised additive models, and explicitly constrained models. Comparing an interpretable model with a more complex alternative would help determine whether post-hoc explanation is necessary for a given task rather than assuming that complexity is always justified.

\subsection{Feature Dependence and Clinical Representation}

Permutation and \gls{shap} analyses should be extended with methods that account more explicitly for correlated predictors. Grouped permutation, conditional importance, or clinically defined feature groups may reduce misleading attribution among substitutable variables. Sensitivity to alternative background samples should be documented.

These extensions would address recognised limitations of marginal permutation and post-hoc attribution methods when predictors are dependent or substitutable \citep{fisher_models_2019, strobl_conditional_2008, vilone_classification_2021}.

Encoded categorical levels and interactions should be mapped back to readable clinical concepts where mathematically defensible. For \gls{hbaone}, future work could compare design-term explanations with grouped explanations at the level of age, specialty, diagnosis, discharge status, and \gls{hbaone} category. Such aggregation must preserve the interaction structure and should not combine coefficients or attributions without a stated rule.

\subsection{Fairness and Subgroup Analysis}

Future work should evaluate predictive performance, calibration, error rates, and explanation behaviour across clinically and demographically relevant subgroups. This is especially important when sex, age, race, or healthcare-utilisation variables contribute to predictions. Subgroup analysis should account for sample size and uncertainty and should not infer fairness from global feature importance alone \citep{allgaier_how_2023, moons_probastai_2025}.

\subsection{Human-Centred and Clinical Evaluation}

The present computational study should be followed by evaluation with intended users before any claim of clinical usefulness. Clinicians should participate in defining relevant explanation tasks, selecting presentation formats, and identifying situations in which explanations help or hinder decisions. Studies should measure decision quality, appropriate reliance, time burden, comprehension, and responses to incorrect but persuasive explanations \citep{panigutti_co-design_2023, vasconcelos_explanations_2023,
lee_understanding_2023, vasey_decideai_2022}.

Any future explanation-screening or warning mechanism would require prospectively justified criteria and human-factors validation. The current agreement and stability results are not sufficient to define a universal acceptance threshold or a clinical safety gate.

\subsection{Reproducibility and Software Development}

The framework should continue to evolve through versioned releases, continuous integration, expanded tests, schema validation, and automated generation of thesis-ready tables and figures from canonical outputs. Duplicate or obsolete result files should be detected by manifests or checksums to prevent stale analyses from entering reports. Model bundles should retain training provenance and, where feasible, a safe and reproducible specification of the training procedure.

These developments would also support the transparent reporting and open-science expectations described by TRIPOD+\gls{ai} \citep{collins_tripodai_2024}.

Containerised execution could complement the existing dependency locks, particularly for long-term preservation. Future archives should record data provenance, code commit, configuration, environment, generated outputs, and comparison reports as one immutable research object.

\section{Code and Data Availability}

The analysis code, notebooks, configuration files, automated tests, and derived result tables supporting this dissertation are maintained in a version-controlled repository and will be released under the MIT License. Dataset provenance and preparation requirements are documented in the repository, while the original datasets remain subject to the access and reuse terms of their respective public sources. After the final code review and archival deposit, this statement will be updated with the public repository address, release version, and persistent \gls{doi}. No new patient data were collected for this study.

\section{Final Remarks}

Explainability is most useful when it increases scrutiny rather than confidence by default. Across the five experiments, the \gls{xai} analyses revealed shared leading features, disagreement among importance definitions, confident prediction errors, and dataset-dependent ranking persistence under bootstrap refitting. None of these computational findings independently established clinical validity.

The central conclusion is that \gls{xai} should be treated as a reproducible model-audit process rather than as an automatic certificate of trust. Applying complementary explanation methods, documenting their assumptions, evaluating stability where possible, and presenting local explanations alongside observed outcomes and prediction uncertainty provides a more complete account of model behaviour \citep{mohseni_multidisciplinary_2021, vilone_classification_2021}. This account remains one component of the broader external, prospective, and human-centred validation required before medical \gls{ai} can responsibly influence clinical decisions \citep{vasey_decideai_2022, panigutti_co-design_2023}.
